\documentclass[../../main.tex]{subfiles}
\begin{document}
\section{Whitney 逼近定理}\label{section:ch06-Whitney逼近定理}

在本节,我们证明本章开头提到的两个关于用光滑映射逼近连续映射的定理。这两个定理与我们刚证明的嵌入定理一样,都归功于 Hassler Whitney [Whi36]。

我们先从关于把函数光滑逼近到 Euclid 空间的定理开始。第一个定理特别表明:从光滑流形 $M$ 到 $\mathbb{R}^k$ 的任一连续函数都可被光滑函数一致逼近。事实上,我们将证明更强的结果。若 $\delta:M\to\mathbb{R}$ 是正连续函数,我们说两个函数 $F,\widetilde{F}:M\to\mathbb{R}^k$ 是 $\delta$-接近的,如果对所有 $x\in M$ 有
\begin{align*}
|F(x)-\widetilde{F}(x)|<\delta(x).
\end{align*}

\begin{theorem}[Whitney Approximation Theorem for Functions]\label{theorem:ch06-approx-functions}
设 $M$ 是带边或不带边光滑流形, $F:M\to\mathbb{R}^k$ 是连续函数。给定任一正连续函数 $\delta:M\to\mathbb{R}$,存在光滑函数 $\widetilde{F}:M\to\mathbb{R}^k$ 与 $F$ 是 $\delta$-接近的。若 $F$ 在闭子集 $A\subseteq M$ 上光滑,则可取 $\widetilde{F}$ 在 $A$ 上等于 $F$。
\end{theorem}

\begin{proof}
若 $F$ 在闭子集 $A$ 上光滑,则由光滑函数的延拓引理(\cref{lemma:ch02-extension})存在光滑函数 $F_0:M\to\mathbb{R}^k$,在 $A$ 上与 $F$ 一致。令
\begin{align*}
U_0=\{y\in M:|F_0(y)-F(y)|<\delta(y)\}.
\end{align*}
则 $U_0$ 是含 $A$ 的开子集。(若没有这样的集合 $A$,我们就取 $U_0=A=\varnothing$ 且 $F_0\equiv 0$。)

我们将证明:存在可数多个点 $\{x_i\}_{i\geqslant 1}\subseteq M\setminus A$ 与 $x_i$ 在 $M\setminus A$ 中的邻域 $U_i$,使 $\{U_i\}_{i\geqslant 1}$ 是 $M\setminus A$ 的开覆盖,且
\begin{align}\label{align:ch06-close}
|F(y)-F(x_i)|<\delta(y)\quad\text{对所有 }y\in U_i.
\end{align}

为看此,对任一 $x\in M\setminus A$,令 $U_x$ 为含于 $M\setminus A$ 的 $x$ 的邻域,并足够小,使对所有 $y\in U_x$ 有
\begin{align*}
\delta(y)>\tfrac12\delta(x)\quad\text{且}\quad |F(y)-F(x)|<\tfrac12\delta(x).
\end{align*}
(由 $\delta$ 与 $F$ 连续,这样的邻域存在。)于是若 $y\in U_x$,有
\begin{align*}
|F(y)-F(x)|<\tfrac12\delta(x)<\delta(y).
\end{align*}
集族 $\{U_x:x\in M\setminus A\}$ 是 $M\setminus A$ 的开覆盖。选取可数子覆盖 $\{U_{x_i}\}_{i\geqslant 1}$ 并令 $U_i=U_{x_i}$,即得\eqref{align:ch06-close}。

令 $\{\varphi_0,\varphi_i\}$ 为从属于 $M$ 的开覆盖 $\{U_0,U_i\}$ 的光滑单位分解,并定义 $\widetilde{F}:M\to\mathbb{R}^k$ 为
\begin{align*}
\widetilde{F}(y)=\varphi_0(y)F_0(y)+\sum_{i\geqslant 1}\varphi_i(y)F(x_i).
\end{align*}
则 $\widetilde{F}$ 显然光滑,且在 $A$ 上等于 $F$。对任一 $y\in M$,由 $\sum_{i\geqslant 0}\varphi_i\equiv 1$ 蕴涵
\begin{align*}
|\widetilde{F}(y)-F(y)|
&=\left|\varphi_0(y)F_0(y)+\sum_{i\geqslant 1}\varphi_i(y)F(x_i)-\left(\varphi_0(y)+\sum_{i\geqslant 1}\varphi_i(y)\right)F(y)\right|\\
&\leqslant \varphi_0(y)|F_0(y)-F(y)|+\sum_{i\geqslant 1}\varphi_i(y)|F(x_i)-F(y)|\\
&<\varphi_0(y)\delta(y)+\sum_{i\geqslant 1}\varphi_i(y)\delta(y)=\delta(y).
\end{align*}
\end{proof}

\begin{corollary}\label{corollary:ch06-small-positive-function}
若 $M$ 是带边或不带边光滑流形, $\delta:M\to\mathbb{R}$ 是正连续函数,则存在光滑函数 $\varepsilon:M\to\mathbb{R}$,使对所有 $x\in M$ 有 $0<\varepsilon(x)<\delta(x)$。
\end{corollary}

\begin{proof}
用 Whitney 逼近定理构造光滑函数 $\varepsilon:M\to\mathbb{R}$,使对所有 $x\in M$ 有
\begin{align*}
\left|\varepsilon(x)-\tfrac12\delta(x)\right|<\tfrac12\delta(x).
\end{align*}
\end{proof}

\subsection{管状邻域}\label{subsection:ch06-tubular}

我们想找到一种方法把 Whitney 逼近定理应用于在光滑流形之间产生连续映射的光滑逼近。若 $F:N\to M$ 是这样一个映射,则由 Whitney 嵌入定理,我们可把 $M$ 视为某 Euclid 空间 $\mathbb{R}^n$ 的嵌入子流形,并用 $\mathbb{R}^n$ 中的光滑映射逼近 $F$。然而一般而言,该光滑映射的像不会落在 $M$ 中。为修正这一点,我们需要知道存在从 $M$ 的某邻域到 $M$ 上的光滑收缩映射。为此,我们引进管状邻域。

对每个 $x\in\mathbb{R}^n$,切空间 $T_x\mathbb{R}^n$ 典范地等同于 $\mathbb{R}^n$ 本身,而切丛 $T\mathbb{R}^n$ 典范微分同胚于 $\mathbb{R}^n\times\mathbb{R}^n$。借助这一等同,每个切空间 $T_x\mathbb{R}^n$ 继承一个 Euclid 点积。

设 $M\subseteq\mathbb{R}^n$ 是嵌入 $m$ 维子流形。对每个 $x\in M$,我们定义 $M$ 在 $x$ 处的法空间为 $T_x\mathbb{R}^n$ 中由关于 Euclid 点积与 $T_xM$ 正交的所有向量组成的 $(n-m)$ 维子空间 $N_xM$。$M$ 的法丛记为 $NM$,是 $T\mathbb{R}^n\cong\mathbb{R}^n\times\mathbb{R}^n$ 中由与 $M$ 正交的向量组成的子集:
\begin{align*}
NM=\{(x,v)\in\mathbb{R}^n\times\mathbb{R}^n:x\in M,\ v\in N_xM\}.
\end{align*}
有一自然投影 $\pi_{NM}:NM\to M$,定义为 $\pi:T\mathbb{R}^n\to\mathbb{R}^n$ 限制到 $NM$ 上。

\begin{theorem}\label{theorem:ch06-normal-bundle-submanifold}
若 $M\subseteq\mathbb{R}^n$ 是嵌入 $m$ 维子流形,则 $NM$ 是 $T\mathbb{R}^n\cong\mathbb{R}^n\times\mathbb{R}^n$ 的 $n$ 维嵌入子流形。
\end{theorem}

\begin{proof}
令 $x_0$ 为 $M$ 的任一点,并令 $(U,\varphi)$ 为 $M$ 在 $\mathbb{R}^n$ 中以 $x_0$ 为中心的切片坐标卡。记 $\widehat{U}=\varphi(U)\subseteq\mathbb{R}^n$,并把 $\varphi$ 的坐标函数记为 $\bigl(u^1,\dots,u^n\bigr)$,使 $M\cap U$ 是 $u^{m+1}=\cdots=u^n=0$ 的集合。在每一点 $x\in U$,向量
\begin{align*}
E_j|_x=(\mathrm{d}\varphi_x)^{-1}\bigl(\partial/\partial u^j|_{\varphi(x)}\bigr)
\end{align*}
构成 $T_x\mathbb{R}^n$ 的一组基。我们可把每个 $E_j|_x$ 按标准标架展开为
\begin{align*}
E_j|_x=E^i_j(x)\frac{\partial}{\partial x^i}\bigg|_x,
\end{align*}
其中每个 $E^i_j(x)$ 是 $\varphi^{-1}$ 在 $\varphi(x)$ 处取值的偏导数,因而是 $x$ 的光滑函数。

定义光滑函数 $\Phi:U\times\mathbb{R}^n\to\widehat{U}\times\mathbb{R}^n$ 为
\begin{align*}
\Phi(x,v)=\bigl(u^1(x),\dots,u^n(x),\ v\cdot E_1|_x,\dots,\ v\cdot E_n|_x\bigr).
\end{align*}
$\Phi$ 在点 $(x,v)$ 处的全导数为
\begin{align*}
\mathrm{D}\Phi(x,v)=
\begin{pmatrix}
\dfrac{\partial u^i}{\partial x^j}(x) & 0\\
0 & E^i_j(x)
\end{pmatrix},
\end{align*}
它可逆,故 $\Phi$ 是局部微分同胚。若 $\Phi(x,v)=\Phi(x_0,v_0)$,则因 $\varphi$ 单射有 $x=x_0$,而对每个 $i$, $v\cdot E_i|_x=v_0\cdot E_i|_x$ 这一事实蕴涵 $v-v_0$ 与 $\bigl(E_1|_x,\dots,E_n|_x\bigr)$ 张成的空间正交,因而为零。于是 $\Phi$ 单射,故它在 $U\times\mathbb{R}^n$ 上定义了光滑坐标卡。定义蕴涵 $(x,v)\in NM$,当且仅当 $\Phi(x,v)$ 落在切片
\begin{align*}
\{(y,z)\in\mathbb{R}^n\times\mathbb{R}^n:y^{m+1}=\cdots=y^n=0,\ z^1=\cdots=z^m=0\}
\end{align*}
中。于是 $\Phi$ 是 $NM$ 在 $\mathbb{R}^n\times\mathbb{R}^n$ 中的切片坐标卡。
\end{proof}

(我们在第十章将能给这个定理一个更简短的证明;见\cref{corollary:ch10-normal-bundle}。)

把 $NM$ 视为 $\mathbb{R}^n\times\mathbb{R}^n$ 的子流形,我们定义 $E:NM\to\mathbb{R}^n$ 为
\begin{align*}
E(x,v)=x+v.
\end{align*}
这映射把每个法空间 $N_xM$ 仿射地映到过 $x$ 且与 $T_xM$ 正交的仿射子空间上。显然 $E$ 光滑,因为它是加法映射 $\mathbb{R}^n\times\mathbb{R}^n\to\mathbb{R}^n$ 限制到 $NM$ 上。$M$ 的一个\textbf{管状邻域}是 $M$ 在 $\mathbb{R}^n$ 中的一个邻域 $U$,它是 $NM$ 的某形如
\begin{align}\label{align:ch06-tubular-V}
V=\{(x,v)\in NM:|v|<\delta(x)\}
\end{align}
的开子集在 $E$ 下的微分同胚像,其中 $\delta:M\to\mathbb{R}$ 是某正连续函数(\cref{figure:ch06-tubular})。

\begin{figure}[ht]
\centering
\includegraphics[width=0.8\textwidth]{fig-06-6.png}
\caption{管状邻域}\label{figure:ch06-tubular}
\end{figure}

\begin{theorem}[Tubular Neighborhood Theorem]\label{theorem:ch06-tubular}
$\mathbb{R}^n$ 的每个嵌入子流形都有管状邻域。
\end{theorem}

\begin{proof}
令 $M_0\subseteq NM$ 为子集 $\{(x,0):x\in M\}$。我们先证明 $E$ 在 $M_0$ 的某邻域上是局部微分同胚。由反函数定理,只需证明微分 $\mathrm{d}E_{(x,0)}$ 在每一点 $(x,0)\in M_0$ 处双射。这由下面两个事实易得:第一, $E$ 在 $M_0$ 上的限制是显然的微分同胚 $M_0\to M$,故 $\mathrm{d}E_{(x,0)}$ 把子空间 $T_{(x,0)}M_0\subseteq T_{(x,0)}NM$ 同构地映到 $T_xM$。第二, $E$ 在纤维 $N_xM$ 上的限制是仿射映射 $w\mapsto x+w$,故 $\mathrm{d}E_{(x,0)}$ 把 $T_{(x,0)}(N_xM)\subseteq T_{(x,0)}NM$ 同构地映到 $N_xM$。由于 $T_x\mathbb{R}^n=T_xM\oplus N_xM$,这表明 $\mathrm{d}E_{(x,0)}$ 满射,从而由维数原因双射。于是 $E$ 在 $NM$ 中 $(x,0)$ 的某邻域上是微分同胚,我们可取它形如 $V_\delta(x)=\{(x_0,v_0)\in NM:|x-x_0|<\delta,\ |v_0|<\delta\}$ 对某 $\delta>0$。(这用到 $NM$ 嵌入在 $\mathbb{R}^n\times\mathbb{R}^n$ 中,故其拓扑由 Euclid 度量诱导。)

为完成证明,我们需证明存在形如\eqref{align:ch06-tubular-V} 的开子集 $V$,使 $E$ 在其上是整体微分同胚。对每个点 $x\in M$,令 $\rho(x)$ 为所有 $\delta\geqslant 1$ 之上确界,使 $E$ 是从 $V_\delta(x)$ 到其像上的微分同胚。上一段的论证蕴涵 $\rho:M\to\mathbb{R}$ 正。为证它连续,任取 $x,x_0\in M$,并先设 $|x-x_0|<\rho(x)$。由三角不等式,当 $\delta=\rho(x)-|x-x_0|$ 时 $V_\delta(x_0)\subseteq V_{\rho(x)}(x)$(\cref{figure:ch06-rho-continuity}),这蕴涵 $\rho(x_0)\geqslant\rho(x)-|x-x_0|$,即
\begin{align}\label{align:ch06-rho-lip}
\rho(x)-\rho(x_0)\leqslant|x-x_0|.
\end{align}
另一方面,若 $|x-x_0|\geqslant\rho(x)$,则\eqref{align:ch06-rho-lip} 平凡成立。交换 $x$ 与 $x_0$ 的角色得类似不等式,这表明 $|\rho(x)-\rho(x_0)|\leqslant|x-x_0|$,故 $\rho$ 连续。注意 $E$ 在整个集合 $V_{\rho(x)}(x)$ 上是单射,因为该集中任两点 $(x_1,v_1),(x_2,v_2)$ 都在某 $V_\delta(x)$ 中,其中 $\delta<\rho(x)$。

现令
\begin{align*}
V=\left\{(x,v)\in NM:|v|<\tfrac12\rho(x)\right\}.
\end{align*}
我们将证明 $E$ 在 $V$ 上单射。设 $(x,v)$ 与 $(x_0,v_0)$ 是 $V$ 中两点使 $E(x,v)=E(x_0,v_0)$(\cref{figure:ch06-E-injective})。不失一般性设 $\rho(x_0)\geqslant\rho(x)$。由 $x+v=x_0+v_0$ 得
\begin{align*}
|x-x_0|=|v-v_0|\leqslant|v|+|v_0|<\tfrac12\rho(x)+\tfrac12\rho(x_0)\leqslant\rho(x).
\end{align*}
因此 $(x,v)$ 与 $(x_0,v_0)$ 都在 $V_{\rho(x)}(x)$ 中。由于 $E$ 在该集上单射,这蕴涵 $(x,v)=(x_0,v_0)$。

集合 $U=E(V)$ 在 $\mathbb{R}^n$ 中开,因为 $E|_V$ 是局部微分同胚从而是开映射。于是 $E:V\to U$ 是光滑双射且为局部微分同胚,故由\cref{proposition:ch04-local-diffeo-equiv} 知它是微分同胚。因此 $U$ 是 $M$ 的管状邻域。
\end{proof}

\begin{figure}[ht]
\centering
\includegraphics[width=0.8\textwidth]{fig-06-7.png}
\caption{$\rho$ 的连续性}\label{figure:ch06-rho-continuity}
\end{figure}

\begin{figure}[ht]
\centering
\includegraphics[width=0.8\textwidth]{fig-06-8.png}
\caption{$E$ 的单射性}\label{figure:ch06-E-injective}
\end{figure}

管状邻域最有用的特征之一由下一命题给出。拓扑空间 $X$ 到子空间 $M\subseteq X$ 的一个\textbf{收缩映射}是连续映射 $r:X\to M$,使 $r|_M$ 是 $M$ 的恒等映射。

\begin{proposition}\label{proposition:ch06-tubular-retraction}
设 $M\subseteq\mathbb{R}^n$ 是嵌入子流形。若 $U$ 是 $M$ 的任一管状邻域,则存在光滑映射 $r:U\to M$,它既是收缩映射又是光滑淹没。
\end{proposition}

\begin{proof}
令 $NM\subseteq T\mathbb{R}^n$ 为 $M$ 的法丛,并令 $M_0\subseteq NM$ 为集合 $M_0=\{(x,0):x\in M\}$。由管状邻域的定义,存在含 $M_0$ 的开子集 $V\subseteq NM$,使 $E:V\to U$ 是微分同胚。定义 $r:U\to M$ 为 $r=\pi_{NM}\circ E^{-1}$,其中 $\pi_{NM}:NM\to M$ 是自然投影。则 $r$ 由复合知光滑。对 $x\in M$,注意 $E(x,0)=x$,故
\begin{align*}
r(x)=\pi_{NM}\circ E^{-1}(x)=\pi_{NM}(x,0)=x,
\end{align*}
这表明 $r$ 是收缩映射。由于 $\pi_{NM}$ 是光滑淹没且 $E^{-1}$ 是微分同胚,故 $r$ 是光滑淹没。
\end{proof}

\subsection{流形之间映射的光滑逼近}\label{subsection:ch06-manifold-approx}

现在我们可把 Whitney 逼近定理推广到流形之间的映射。这一推广将在第十六章对曲线积分、在第十七至十八章对 de Rham 上同调有重要应用。

\begin{theorem}[Whitney Approximation Theorem]\label{theorem:ch06-approx-maps}
设 $N$ 是带边或不带边光滑流形, $M$ 是光滑流形(不带边), $F:N\to M$ 是连续映射。则 $F$ 同伦于某光滑映射。若 $F$ 已在闭子集 $A\subseteq N$ 上光滑,则同伦可取得关于 $A$ 相对。
\end{theorem}

\begin{proof}
由 Whitney 嵌入定理,我们不妨设 $M$ 是 $\mathbb{R}^n$ 的真嵌入子流形。令 $U$ 为 $M$ 在 $\mathbb{R}^n$ 中的管状邻域,并令 $r:U\to M$ 为\cref{proposition:ch06-tubular-retraction} 给出的光滑收缩映射。对任一 $x\in M$,令
\begin{align}\label{align:ch06-delta-tube}
\delta(x)=\sup\{\varepsilon\geqslant 1:B_\varepsilon(x)\subseteq U\}.
\end{align}
由一个与管状邻域定理证明中类似的三角不等式论证, $\delta:M\to\mathbb{R}^+$ 连续。令 $\widetilde{\delta}=\delta\circ F:N\to\mathbb{R}^+$。由\cref{theorem:ch06-approx-functions},存在光滑函数 $\widetilde{F}:N\to\mathbb{R}^n$ 与 $F$ 是 $\widetilde{\delta}$-接近的,且在 $A$ 上等于 $F$($A$ 可以是空集)。令 $H:N\times I\to M$ 为 $r$ 与 $F,\widetilde{F}$ 之间直线同伦的复合:
\begin{align*}
H(p,t)=r\bigl((1-t)F(p)+t\widetilde{F}(p)\bigr).
\end{align*}
这是良定义的,因为我们对 $\widetilde{F}$ 的条件保证:对每个 $p\in N$,
\begin{align*}
|\widetilde{F}(p)-F(p)|<\widetilde{\delta}(p)=\delta(F(p)),
\end{align*}
这意味着 $\widetilde{F}(p)$ 含于以 $F(p)$ 为中心、半径 $\delta(F(p))$ 的球中;由于该球含于 $U$,故从 $F(p)$ 到 $\widetilde{F}(p)$ 的整条线段也在 $U$ 中。

于是 $H$ 是 $H(p,0)=F(p)$ 与 $H(p,1)=r\bigl(\widetilde{F}(p)\bigr)$ 之间的同伦,后者由复合知是光滑映射。它满足对所有 $p\in A$ 有 $H(p,t)=F(p)$,因为在那里 $F=\widetilde{F}$。
\end{proof}

\begin{corollary}[Extension Lemma for Smooth Maps]\label{corollary:ch06-extension-maps}
设 $N$ 是带边或不带边光滑流形, $M$ 是光滑流形, $A\subseteq N$ 是闭子集, $f:A\to M$ 是光滑映射。则 $f$ 有到 $N$ 的光滑延拓,当且仅当它有到 $N$ 的连续延拓。
\end{corollary}

\begin{proof}
若 $F:N\to M$ 是 $f$ 到整个 $N$ 的连续延拓, Whitney 逼近定理保证存在光滑映射 $\widetilde{F}$(同伦于 $F$,事实上如此,尽管我们此处不需要)在 $A$ 上与 $f$ 一致;换言之, $\widetilde{F}$ 是 $f$ 的光滑延拓。逆命题显然。
\end{proof}

若 $N$ 与 $M$ 是两个带边或不带边光滑流形,同伦 $H:N\times I\to M$ 称为\textbf{光滑同伦},如果它也是光滑映射,即它可延拓为 $N\times\mathbb{R}$ 中 $N\times I$ 某邻域上的光滑映射。两个映射称为光滑同伦的,如果它们之间存在光滑同伦。

\begin{lemma}\label{lemma:ch06-smooth-homotopy-eq}
若 $N$ 与 $M$ 是带边或不带边光滑流形,则光滑同伦在所有从 $N$ 到 $M$ 的光滑映射之集上是等价关系。
\end{lemma}

\begin{proof}
自反性与对称性的证明与普通同伦相同。为证传递性,设 $F,G,K:N\to M$ 为光滑映射, $H_1,H_2:N\times I\to M$ 分别为从 $F$ 到 $G$ 与从 $G$ 到 $K$ 的光滑同伦。令 $\psi:[0,1]\to[0,2]$ 为光滑映射,使
\begin{align*}
0\leqslant\psi(t)\leqslant 1\text{ 对 }t\in\left[0,\tfrac12\right],\quad
1\leqslant\psi(t)\leqslant 2\text{ 对 }t\in\left[\tfrac12,1\right],
\end{align*}
$\psi(0)=0$, $\psi(1)=2$,且 $\psi(t)\equiv 1$ 对 $t$ 在 $\tfrac12$ 的某邻域内(见\cref{figure:ch06-psi})。定义 $H:N\times I\to M$ 为
\begin{align*}
H(x,t)=
\begin{cases}
H_1\bigl(x,\psi(t)\bigr), & t\in\left[0,\tfrac12\right],\\
H_2\bigl(x,\psi(t)-1\bigr), & t\in\left[\tfrac12,1\right].
\end{cases}
\end{align*}
于是易验证 $H$ 是从 $F$ 到 $K$ 的光滑同伦。
\end{proof}

\begin{figure}[ht]
\centering
\includegraphics[width=0.8\textwidth]{fig-06-9.png}
\caption{引理 6.28 证明中的函数 $\varphi$}\label{figure:ch06-psi}
\end{figure}

\begin{theorem}\label{theorem:ch06-homotopic-smoothly}
设 $N$ 是带边或不带边光滑流形, $M$ 是光滑流形, $F,G:N\to M$ 是光滑映射。若 $F$ 与 $G$ 同伦,则它们光滑同伦。若 $F$ 与 $G$ 关于某闭子集 $A\subseteq N$ 同伦,则它们关于 $A$ 光滑同伦。
\end{theorem}

\begin{proof}
设 $F,G:N\to M$ 光滑,并令 $H:N\times I\to M$ 为从 $F$ 到 $G$ 的同伦(关于 $A$, $A$ 可以是空集)。我们要证明 $H$ 可被光滑同伦替代。

定义 $\widetilde{H}:N\times\mathbb{R}\to M$ 为
\begin{align*}
\widetilde{H}(x,t)=
\begin{cases}
H(x,t), & t\in[0,1],\\
H(x,0), & t\leqslant 0,\\
H(x,1), & t\geqslant 1.
\end{cases}
\end{align*}
由粘合引理,它连续。$\widetilde{H}$ 在 $N\times\{0\}\cup N\times\{1\}$ 上的限制光滑,因为它在 $N\times\{0\}$ 上等于 $F\circ\pi_1$,在 $N\times\{1\}$ 上等于 $G\circ\pi_1$(其中 $\pi_1:N\times\mathbb{R}\to N$ 是投影)。若 $H$ 是关于 $A$ 的同伦,则 $\widetilde{H}$ 在 $A\times I$ 上也光滑。由于 $N\times\mathbb{R}$ 是带(可能空)边界的光滑流形, Whitney 逼近定理保证存在光滑映射 $\widehat{H}:N\times\mathbb{R}\to M$(同伦于 $\widetilde{H}$,但此处不需要),其在 $N\times\{0\}\cup N\times\{1\}\cup A\times I$ 上的限制与 $\widetilde{H}$(从而与 $H$)一致。再次限制回 $N\times I$,我们看到 $\widehat{H}|_{N\times I}$ 是 $F$ 与 $G$ 之间的光滑同伦(关于 $A$)。
\end{proof}

当目标流形具有非空边界时, \cref{theorem:ch06-approx-maps} 与\cref{theorem:ch06-homotopic-smoothly} 的类似结论不成立,因为可能找不到在 $A$ 上与 $F$ 一致的光滑映射(见习题 6-7)。然而,若我们不坚持关于某子集相对的同伦,其余结果可推广到到带边流形的映射。证明将推迟到第九章(见\cref{theorem:ch09-approx-boundary} 与\cref{theorem:ch09-homotopy-boundary})。

\end{document}